% GENERATED FILE. Edit canonical lessons, then run npm run generate:books.
% canonical-source-hash: fnv1a64:95982a1ab19d6c98
% canonical-lessons: GE-C06-neun, GE-C06-zehn, GE-C06-zahlen-6-10, GE-C06-monate-latein, GE-C06-practice

\chapter{Nine, Ten, and the Months That Are Numbers}
\label{ch:ge-zahlen-9-10}
\hlchaptermodality{9}

\begin{chapteropening}
\textbf{By the end of this chapter:} \emph{I can count all the way to ten in German, and say which words a language keeps and which it buys.}

Nine and ten finish the count, and the chapter closes on the seam the whole sequence has been building toward: the numbers are homegrown and the month names are Latin imports, so one calendar carries both answers to the same question.
\end{chapteropening}

\section[neun]{neun — \textquotedblleft{}nine,\textquotedblright{} in the vowel of \emph{Deutsch}}

\label{lesson:GE-C06-neun}



\begin{quote}
Ninth number, and again nothing new to learn in the mouth — you got this vowel from the name of the language itself.
\end{quote}

\subsection*{You'll want to know: neun}
\begin{quote}
\textbf{neun} — \textquotedblleft{}nine.\textquotedblright{} Said \emph{noyn}.
\end{quote}

The \textbf{eu} is the \emph{oy} of English \emph{boy}, which is exactly what it does in \emph{Deutsch}. So \emph{neun} rhymes with English \emph{coin}, and it is one syllable with an \emph{n} at each end.

Do not let the spelling talk you into \emph{new-n}. German \textbf{eu} is never two separate vowels.

\begin{cousinweb}
\textbf{neun} is English \textbf{nine}, and the two have kept the same skeleton: \emph{n}, a vowel, \emph{n}. Only the vowel wandered — German to \emph{oy}, English to \emph{eye}.

The Latin cousin is the odd one out here. \emph{Novem} has a \emph{v} where the Germanic words have nothing, which is why Spanish \emph{nueve} and French \emph{neuf} look like a different word. Between \emph{neun} and \emph{nine} there is nothing to explain.
\end{cousinweb}

\subsection*{Your turn}
Say these aloud:

\begin{itemize}
\raggedright
  \item \textquotedblleft{}neun\textquotedblright{} — \emph{noyn}, rhyming with \emph{coin}
  \item \textquotedblleft{}Deutsch\textquotedblright{} then \textquotedblleft{}neun\textquotedblright{} — the same \emph{eu}
  \item \textquotedblleft{}sieben, acht, neun\textquotedblright{}
\end{itemize}

\begin{tcolorbox}[breakable,colback=teal!4,colframe=teal!35!black,title={Before you move on}]
What is \textquotedblleft{}nine\textquotedblright{}? (\emph{Neun}, said \emph{noyn}.) How is \textbf{eu} said? (Like the \emph{oy} of \emph{boy}, as in \emph{Deutsch}.) What do \emph{neun} and \emph{nine} share? (\emph{N}, a vowel, \emph{n} — only the vowel moved.)
\end{tcolorbox}
\section[zehn]{zehn — \textquotedblleft{}ten,\textquotedblright{} and one consonant in three photographs}

\label{lesson:GE-C06-zehn}



\begin{quote}
The last one. And it is the number that shows both of German's sound laws working, one after the other, on the same letter.
\end{quote}

\subsection*{You'll want to know: zehn}
\begin{quote}
\textbf{zehn} — \textquotedblleft{}ten.\textquotedblright{} Said \emph{tsayn}.
\end{quote}

The \textbf{z} is the \emph{ts} you met in \emph{zwei}, and the \textbf{h} does what it did in \emph{wohnen} and \emph{gehen}: it says nothing and stretches the vowel in front of it. So: \emph{ts}, a long \emph{ay}, and an \emph{n}.

\begin{cousinweb}
Latin \emph{decem}, English \textbf{ten}, German \textbf{zehn}. One word, photographed at three moments of the same consonant walking:

\noindent
\begin{tabularx}{\linewidth}{@{}>{\raggedright\arraybackslash}X>{\raggedright\arraybackslash}X>{\raggedright\arraybackslash}X@{}}
\toprule
\textbf{Stage} & \textbf{Word} & \textbf{The sound} \\
\midrule
what Latin kept & \emph{decem} & \textbf{d} \\
the first law & \textbf{ten} & \textbf{d} became \textbf{t} \\
the second law & \emph{zehn} & \textbf{t} became \textbf{ts} \\
\bottomrule
\end{tabularx}

The first step is the law from the \emph{drei} lesson, the one that separated Germanic from Latin. The second is the \emph{later} shift you already know from \emph{gut} and \emph{machen} — the one that only German went through, which is why English stopped at \emph{ten} and German carried on to \emph{zehn}.

So the distance between \emph{ten} and \emph{zehn} is not a different word. It is one extra step down a road English stopped walking.
\end{cousinweb}

\subsection*{Your turn}
Say these aloud:

\begin{itemize}
\raggedright
  \item \textquotedblleft{}zehn\textquotedblright{} — \emph{tsayn}, long vowel, silent \emph{h}
  \item \textquotedblleft{}decem, ten, zehn\textquotedblright{} — the same consonant, three times
  \item \textquotedblleft{}acht, neun, zehn\textquotedblright{}
\end{itemize}

\begin{tcolorbox}[breakable,colback=teal!4,colframe=teal!35!black,title={Before you move on}]
What is \textquotedblleft{}ten\textquotedblright{}? (\emph{Zehn}, said \emph{tsayn}.) What does the \textbf{h} do? (Nothing — it lengthens the vowel.) Walk \emph{decem} to \emph{zehn}. (\emph{D} became \emph{t}, and German took the \emph{t} on to \emph{ts}.)
\end{tcolorbox}
\section[Practice]{sechs bis zehn — finishing the count}

\label{lesson:GE-C06-zahlen-6-10}



\begin{quote}
Five more words learned one at a time. Same job as before: make them one run, then join the two runs together.
\end{quote}

\subsection*{You'll want to know: the second run}
\begin{quote}
\textbf{sechs, sieben, acht, neun, zehn}
\end{quote}

Said straight through: \emph{zeks, ZEE-ben, akht, noyn, tsayn}.

And the second five are the same story as the first — English twins, Latin only a cousin:

\noindent
\begin{tabularx}{\linewidth}{@{}>{\raggedright\arraybackslash}X>{\raggedright\arraybackslash}X>{\raggedright\arraybackslash}X@{}}
\toprule
\textbf{German} & \textbf{English} & \textbf{What changed} \\
\midrule
\emph{sechs} & six & almost nothing; \emph{chs} is \emph{x} \\
\emph{sieben} & seven & the \emph{v} softened to a \emph{b} \\
\emph{acht} & eight & German still says the scrape English writes as \textbf{gh} \\
\emph{neun} & nine & only the vowel moved \\
\emph{zehn} & ten & an old \textbf{d} to \textbf{t}, then \textbf{t} to \textbf{ts} \\
\bottomrule
\end{tabularx}

\subsection*{Your turn}
Say these aloud:

\begin{itemize}
\raggedright
  \item \textquotedblleft{}sechs, sieben, acht, neun, zehn\textquotedblright{} — straight through
  \item all ten from the start, in one breath if you can
  \item the ten backwards, from \textquotedblleft{}zehn\textquotedblright{} down to \textquotedblleft{}eins\textquotedblright{}
\end{itemize}

\emph{Twice through:} \textquotedblleft{}eins, zwei, drei, vier, fünf, sechs, sieben, acht, neun, zehn\textquotedblright{}

\begin{tcolorbox}[breakable,colback=teal!4,colframe=teal!35!black,title={Before you move on}]
Count from six to ten. (\emph{Sechs, sieben, acht, neun, zehn}.) Now all ten. Which of the ten is borrowed from Latin? (None — every one is the English number's twin.)
\end{tcolorbox}
\section[(homegrown numbers, imported months)]{Homegrown numbers, imported months}

\label{lesson:GE-C06-monate-latein}



\begin{quote}
You have ten numbers and not one of them came from Latin. That is worth holding onto, because the very next thing German counts with \textbf{did} come from Latin — and you can hear the join.
\end{quote}

\begin{grammarlens}[title={Grammar Lens: a language keeps some words and buys others}]
A language does not import evenly. The words people use every hour of every day are the hardest to displace, and the words that arrive with an institution come in whole.

Numbers are the first kind. German counts with its own inherited words, and so does English. Month names are the second kind: they arrived with the Roman calendar, and both languages simply took them.
\end{grammarlens}

\begin{culture}
The Roman year began in \textbf{March}. Count from there and the last four months of the year land on the Latin words for seven, eight, nine and ten — which is why those four month names are numbers wearing a disguise, in German exactly as in English.

Say each Latin number and then the German one that means the same thing:

\noindent
\begin{tabularx}{\linewidth}{@{}>{\raggedright\arraybackslash}X>{\raggedright\arraybackslash}X@{}}
\toprule
\textbf{The Latin number inside} & \textbf{German counts it} \\
\midrule
\emph{septem} — seven & \emph{sieben} \\
\emph{octo} — eight & \emph{acht} \\
\emph{novem} — nine & \emph{neun} \\
\emph{decem} — ten & \emph{zehn} \\
\bottomrule
\end{tabularx}

Not one pair sounds like a pair. The month names came from one language and the counting words from another, and German never reconciled them — because there was never any reason to.

You will meet the German month names themselves in a later chapter. What matters now is that when you get there, they will look Latin, and the numbers you have just learned will still be German.
\end{culture}

\subsection*{Your turn}
Say these aloud:

\begin{itemize}
\raggedright
  \item \textquotedblleft{}sieben, acht, neun, zehn\textquotedblright{} — the German four
  \item \textquotedblleft{}septem, octo, novem, decem\textquotedblright{} — the Latin four
  \item which kind of word a language keeps, and which it buys
\end{itemize}

\begin{tcolorbox}[breakable,colback=teal!4,colframe=teal!35!black,title={Before you move on}]
Did German borrow its numbers? (No — they are its own.) Did it borrow its month names? (Yes, from Latin, with the calendar.) Why do four of them carry the wrong number? (The Roman year began in March, so \emph{septem} was the seventh month.)
\end{tcolorbox}
\section[Practice]{Practice — counting, and asking for one}

\label{lesson:GE-C06-practice}



\begin{quote}
Ten numbers, one article, and one fact about which words a language keeps. Put them to work.
\end{quote}

\begin{grammarlens}[title={Grammar Lens: the ten, and what each one cost}]
\noindent
\begin{tabularx}{\linewidth}{@{}>{\raggedright\arraybackslash}X>{\raggedright\arraybackslash}X>{\raggedright\arraybackslash}X@{}}
\toprule
\textbf{German} & \textbf{Said} & \textbf{English twin} \\
\midrule
\emph{eins} & \emph{eyns} & one \\
\emph{zwei} & \emph{tsvy} & two \\
\emph{drei} & \emph{dry} & three \\
\emph{vier} & \emph{feer} & four \\
\emph{fünf} & \emph{fewnf} & five \\
\emph{sechs} & \emph{zeks} & six \\
\emph{sieben} & \emph{ZEE-ben} & seven \\
\emph{acht} & \emph{akht} & eight \\
\emph{neun} & \emph{noyn} & nine \\
\emph{zehn} & \emph{tsayn} & ten \\
\bottomrule
\end{tabularx}

Ten numbers and not one of them borrowed. Every one is the English word with a predictable coat on, which is the whole reason this chapter was cheap.
\end{grammarlens}

\subsection*{The exchange}
\noindent
\begin{tabularx}{\linewidth}{@{}>{\raggedright\arraybackslash}X>{\raggedright\arraybackslash}X@{}}
\toprule
\textbf{German} & \textbf{English} \\
\midrule
\emph{Eins, zwei, drei …} & One, two, three… \\
\emph{Und?} & And? \\
\emph{… vier, fünf, sechs, sieben, acht, neun, zehn!} & …four, five, six, seven, eight, nine, ten! \\
\emph{Sehr gut!} & Very good! \\
\emph{Danke!} — \emph{Bitte.} & Thanks! — You're welcome. \\
\bottomrule
\end{tabularx}

Every word in that exchange is one you already own, and seven of them you learned this chapter.

Then the article, on the two nouns \emph{guten Tag} and \emph{gute Nacht} are built from:

\begin{quote}
\emph{der Tag} $\to$ \textbf{ein Tag} · \emph{die Nacht} $\to$ \textbf{eine Nacht}
\end{quote}

\emph{Eins} on its own keeps its \textbf{-s}, because nothing follows it. The moment a naming word arrives, the \emph{-s} comes off.

\subsection*{Your turn}
Say these aloud:

\begin{itemize}
\raggedright
  \item all ten, straight through
  \item the exchange, both voices
  \item \textquotedblleft{}eins\textquotedblright{} as a count, then \textquotedblleft{}ein\textquotedblright{} as an article — the \emph{-s} comes off
\end{itemize}

\emph{Twice through:} \textquotedblleft{}eins, zwei, drei, vier, fünf, sechs, sieben, acht, neun, zehn\textquotedblright{}

\begin{tcolorbox}[breakable,colback=teal!4,colframe=teal!35!black,title={Before you move on}]
Count to ten. When does \emph{eins} become \emph{ein}? (In front of a naming word.) Are the German numbers borrowed from Latin? (No — they are the English numbers' twins.) Which German words \emph{were} borrowed from Latin? (The month names.)
\end{tcolorbox}
